\section{Instâncias e Metodologia}\label{sec:instancias}

\subsection{Escolha das instâncias}

Usamos 12 instâncias, em três grupos (Tabela~\ref{tab:familias}). O grupo A tem grafos aleatórios esparsos ($m=2n$, grau médio 4, como a R100k do Trabalho~1), pequenos o bastante para o ótimo ser calculado. O grupo B tem grades toroidais com pesos $\pm1$ e lado curto $L$ pequeno: a largura da programação dinâmica é $2L+1$ e não depende do comprimento da grade, então o ótimo sai até com 100 mil vértices. O grupo C tem instâncias em que a programação dinâmica é inviável (largura de 41 ou mais): ali só roda o Branch-and-Bound, que devolve a melhor solução achada em 15 minutos. Duas delas, o G32 e o G57, vêm do Trabalho~1.

Não usamos o G-set inteiro porque, para essas instâncias, a melhor ordem que o guloso encontra tem fronteira de 41 no G32 e acima de 60 nas demais (a tabela teria pelo menos $2^{41}$ entradas, contra o limite prático de $2^{26}$). Isso não quer dizer que elas não possam ser resolvidas, só que estão fora do alcance do que implementamos.

\begin{table}[ht]\centering\footnotesize
\caption{Instâncias do experimento. $w$: menor largura de fronteira que o guloso encontrou (``$>$'': desistiu acima do valor indicado).}\label{tab:familias}
\begin{tabular}{llrcl}\hline
Grupo & Instância & $n$ & $w$ & Descrição\\\hline
A & \texttt{01\_ALE\_n100} & 100 & 21 & aleatória esparsa, pesos $+1$\\
A & \texttt{02\_ALE\_n120} & 120 & 26 & aleatória esparsa, pesos $+1$\\
A & \texttt{03\_PM1\_n120} & 120 & 26 & aleatória esparsa, pesos $\pm1$\\
B & \texttt{04\_TORO\_8x250} & 2\,000 & 17 & grade toroidal $8\times250$, pesos $\pm1$\\
B & \texttt{05\_TORO\_10x200} & 2\,000 & 21 & grade toroidal $10\times200$, pesos $\pm1$\\
B & \texttt{06\_TORO\_5x20000} & 100\,000 & 11 & grade toroidal $5\times20000$, pesos $\pm1$\\
B & \texttt{07\_TORO\_7x14286} & 100\,002 & 15 & grade toroidal $7\times14286$, pesos $\pm1$\\
C & \texttt{08\_ALE\_n200} & 200 & 42 & aleatória esparsa, pesos $+1$\\
C & \texttt{09\_PM1\_n500} & 500 & $>80$ & aleatória esparsa, pesos $\pm1$\\
C & \texttt{10\_TORO\_20x50} & 1\,000 & 41 & grade toroidal $20\times50$, pesos $\pm1$\\
C & \texttt{11\_G32} & 2\,000 & 41 & G-set (Trabalho 1), pesos $\pm1$\\
C & \texttt{12\_G57} & 5\,000 & $>60$ & G-set (Trabalho 1), pesos $\pm1$\\\hline
\end{tabular}
\end{table}

\subsection{Geradores}

Usamos o gerador do Trabalho~1 e um novo, \texttt{gerador\_t2.cpp}, com o mesmo formato de arquivo (\texttt{n m} e uma linha \texttt{u v peso} por aresta). Ele tem dois modos, ambos determinísticos para uma dada semente: \texttt{aleatorio n m semente pesos}, que sorteia $m$ arestas distintas (pesos todos $+1$ ou $\pm1$ com probabilidade $1/2$); e \texttt{toro L C semente}, que gera uma grade toroidal $L\times C$ (cada vértice com 4 vizinhos), com pesos $\pm1$ e rótulos embaralhados para que a estrutura de grade não apareça na numeração. Todas as instâncias geradas usam a semente 1.

\subsection{Plataforma e protocolo}

Os experimentos rodaram em um computador com processador Intel Core i5-9300H a 2,40\,GHz e 8\,GB de RAM, em execução \emph{single-thread}, sob Windows 11, com o compilador \texttt{g++} 15.1 e as opções \texttt{-std=c++17 -O2}. Os tempos são de relógio de parede e medem só o algoritmo, sem a leitura do arquivo. Cada instância foi executada uma única vez, com limite de 900\,s (15 minutos) \emph{por algoritmo}. Se o Branch-and-Bound estoura o limite, devolve a melhor solução encontrada e a otimalidade não é provada. A programação dinâmica roda com $w\le26$ e até 2\,GB de memória; ela não devolve solução parcial: ou termina com o ótimo, ou não executa. O Branch-and-Bound não roda para $n>5\,000$, o que afeta apenas as instâncias 06 e 07.

A heurística do Trabalho~1 está incorporada ao programa novo, mas a leitura agora ordena as listas de vizinhos, o que muda o desempate da ordenação por grau. É o mesmo algoritmo, mas não idêntico bit a bit: no G32 ele dá 1\,170 aqui contra 1\,154 no relatório anterior, e no G57, 2\,924 contra 2\,848. Neste relatório usamos sempre a versão incorporada.

\subsection{Validação}

Um programa de validação (\texttt{validacao.cpp}) gera 3\,000 grafos aleatórios com 1 a 16 vértices, densidades de 0 a 1 e quatro regimes de pesos (todos $+1$; $\pm1$; inteiros em $[-4,6]$, inclusive nulos; inteiros em $[1,9]$) e compara o ótimo obtido por força bruta com o dos dois algoritmos, recalculando também o corte das partições devolvidas. A força bruta é usada só como oráculo de teste. Essa validação foi feita durante o desenvolvimento, em um ambiente Linux com \texttt{g++} 13.3: 3\,000 casos, 0 falhas. Além disso, os ótimos das instâncias 01 a 07 obtidos nesse ambiente coincidem exatamente com os da máquina do experimento, e a partição final de cada execução é sempre reavaliada pelo programa.

\subsection{Reprodução}

O pacote entregue contém código, instâncias e scripts. A partir da pasta do projeto:
\begin{verbatim}
bash scripts/rodar_experimentos_longos.sh 900
python3 scripts/gerar_tabelas_figuras.py
./exato instancias_longas/04_TORO_8x250
\end{verbatim}
O primeiro comando roda as 12 instâncias (15 minutos por algoritmo), o segundo gera tabelas e figuras a partir do CSV, e o terceiro roda uma instância (\texttt{./exato} sem argumentos abre um menu). As instâncias estão em \texttt{instancias\_longas/} e podem ser recriadas com \texttt{scripts/gerar\_instancias\_longas.sh}.
